\subsection{Hashing}

Permite comparar strings em $O(1)$. Nunca é 100\% deterministicamente correto. Sempre há risco de colisão.
$$\text{hash}(s) = \sum\limits_{i=0}^{n-1} s[i] \cdot p^i \mod m$$
onde $p, m$ são escolhidos. Essa é uma polynomial rolling hashing function. Geralmente escolhe-se $p$ como um número perto do número de letras do alfabeto por ex $p = 31$. Se também tem uppercase usa $p=53$. Pra minimizar colisão $m$ deve ser muito pequeno. Na prática $m=2^{64}$ não é recomendado. Uma boa escolha é um número primo grande por ex $m = 10^9 + 9$. Use a conversão $a \to 1$, $b \to 2$, ..., $z \to 26$.

\inccode{strings/compute_hash.cpp}

Precomputar potências de $p$ pode ser uma boa ideia.

\prob{1}{hash} Encontre strings duplicadas em array de strings e divida-as em grupos. 

\textbf{Solução:} Com sort normal seria $O(nm \log n)$ mas com hash fica sendo $O(nm + n \log n)$.

\inccode{strings/group_identical_strings.cpp}

\prob{2}{hash} Calcule o hash de $s[i..j])$

\textbf{Solução:} $$\text{hash}(s[i..j]) = (\text{hash}(s[0..j]) - \text{hash}(s[0..i-1])) \cdot p^{-i} \mod m$$
podemos precomputar os inversos pra ficar em $O(1)$. No entanto, ter os hashs multiplicados por $p^i$ já é suficiente pra compará-los porque é só multiplicar o de menor potência por $p^{j-i}$.

\prob{3}{hash} Determine o número de substrings diferentes numa string

\textbf{Solução:} Iteramos por todas as substrings. PS: isso demora $O(n^2)$ então use outra coisa se quiser mesmo fazer isso!
\inccode{strings/number_diff_strings.cpp}

Para diminuir a probabilidade de colisão podemos utilizar dois hashings com módulos diferentes (por meio de um pair). Por exemplo $m=10^9+9$ e $m=10^9+7$

\subsubsection{Rabin-Karp para matching}

Dadas duas strings: um texto $t$ e um padrão $s$, determinar se o padrão aparece no texto e se aparecer, enumere todas as ocorrências em $O(|s|+|t|)$.

\inccode{strings/rabinkarp.cpp}